\section{Implementation}
\label{sec:implementation}

\subsection{Overview of the Statutory Mechanism}

The \dia{} has six implementation components:
\begin{enumerate}
  \item The \ar{} default run, implementing the \hh{} bisection tree.
  \item The content-derived seed formula.
  \item The county-preservation default ($\alpha_\text{county} = 2.0$).
  \item The VRA overlay for Gingles-qualified states.
  \item The \callais{} preflight gate enforcing signal disentanglement.
  \item The gcd-incompatibility fallback for states where the seat count
    has a factorization incompatible with a strict \hh{} tree.
\end{enumerate}

We address each in turn.

\subsection{Component 1: The ApportionRegions Default Run}

The default redistricting run implements the \hh{} bisection tree on the
state's TIGER/Line adjacency graph.

The algorithm proceeds as follows.
Let $k$ be the state's apportioned seat count and $k = p_1^{a_1} \cdot p_2^{a_2} \cdots p_r^{a_r}$
be its prime factorization, with factors in non-increasing order.
\begin{enumerate}
  \item Compute the \hh{} bisection tree from $k$'s prime factorization.
    At each level of the tree, the split is into $p_i$ equal sub-groups.
  \item Initialize METIS with the state's TIGER adjacency graph, with census
    tract populations as vertex weights and shared boundary lengths as
    edge weights \citep{b2paper}.
  \item For each node in the bisection tree, from root to leaves:
    \begin{enumerate}
      \item Identify the sub-graph corresponding to the current geographic region.
      \item Run METIS $k$-way partitioning, where $k = p_i$ is the split factor
        at this level.
      \item Apply the \hh{} priority rule to assign census tracts to the $p_i$
        child regions.
    \end{enumerate}
  \item Output the resulting $k$ districts, each consisting of a set of
    census tracts.
\end{enumerate}

The population balance tolerance is set to $\texttt{ubvec}[0] = 1.001$,
imposing a maximum $\pm 0.05\%$ deviation from the ideal district population
$P/k$.
This satisfies the \karcher{} population equality standard with a large margin.

The algorithm is implemented in the \texttt{redist} Rust binary
\citep{b0paper}, which runs a 50-state apportionment in approximately
30 minutes on a laptop.
The \dia{} would require states to run the algorithm within 30 days of
the release of the decennial census data, with maps submitted to the
Election Assistance Commission for verification.

\subsection{Component 2: The Content-Derived Seed Formula}

The random seed for METIS initialization is computed as:
\[
  \texttt{seed} = \text{SHA-256}(\texttt{census\_release\_id}\ \|\ \texttt{`DIA\_SEED\_V1'})
\]
where \texttt{census\_release\_id} is the Census Bureau's official public
identifier for the decennial redistricting data release (e.g.,
\texttt{P.L.\ 94-171\ 2020-08-12}), and \texttt{`DIA\_SEED\_V1'} is a
fixed string identifying the statute version.

The formula has three properties:
\begin{itemize}
  \item \textbf{Deterministic}: Given the census release identifier, the seed
    is uniquely determined.
    Anyone can compute it.
  \item \textbf{Content-derived}: The seed depends on the census data itself,
    not on any choice made by a political actor.
    The Census Bureau's release identifier is fixed at the moment the data
    is released, before any state has begun redistricting.
  \item \textbf{Audit-ready}: The formula is specified in the statute.
    Any court, expert, or citizen can verify that the correct seed was used
    by computing SHA-256 of the census release identifier and comparing it
    to the seed recorded in the plan manifest.
\end{itemize}

Because a single seed might land in a suboptimal local minimum, the statute
does not stop at the content-derived seed.
Instead, the algorithm runs forward from that seed — seed$_0$, seed$_0$+1,
seed$_0$+2, \ldots --- until \textbf{600 consecutive seeds} produce no
improvement in normalised edge cut.
This \emph{Convergence Sweep} (B.16, \citealt{b16paper}) certifies that the
returned map is the empirical minimum: across all 50 states at the 2020
census, the worst-case convergence tail was 511 consecutive non-improving seeds
(Georgia).
Setting $T_{\text{stat}} = 600$ gives an 89-seed margin above this worst case.

The formula therefore serves two roles simultaneously:
(1) a canonical audit reference for the starting seed, and
(2) the anchor of a deterministic convergence walk that terminates with
a certified-optimal map.
\begin{quote}
\textit{Statutory text (2 U.S.C.\ §2a(c)(2)):}
``The algorithm shall iterate over seeds beginning at the content-derived
seed value, halting when no improvement in normalised edge cut has been
observed for 600 consecutive seeds, and shall use the minimum edge cut
plan found across all seeds evaluated.''
\end{quote}
\noindent The combined guarantee: any party can verify the starting seed,
and the stopping criterion is both specified in statute and proven sufficient
by the B.16 analysis of US census tract graphs.

\subsection{Component 3: County Preservation by Default}

The \dia{} specifies a default county-preservation weight of $\alpha_\text{county} = 2.0$,
implemented via the B.10 subdivision-respecting edge-weighting scheme \citep{b10paper}.

At $\alpha_\text{county} = 2.0$, county boundaries are weighted twice as
heavily as other geographic boundaries in the edge-cut criterion.
The algorithm therefore strongly prefers to keep counties intact within a
single district.
When a county must be split (because its population exceeds the ideal district
population), the algorithm minimizes the number of county splits.

The county-preservation default is a statutory default, not an absolute
requirement.
The proposed statute text specifies $\alpha_\text{county} = 2.0$ as the
default, but provides that a state may petition the Election Assistance
Commission for a different weight, subject to demonstrating that the
alternative weight is necessary to comply with the VRA or with the
state's constitution.

This provision reflects the reality that county-preservation is a legitimate
redistricting criterion---embodied in North Carolina's county-grouping
requirement upheld in \textit{Stephenson v.\ Bartlett}---but that
over-weighting county boundaries can conflict with population equality
in states with large counties.
Texas's Harris County (Houston), with a population of 4.7 million, would
require five congressional districts regardless of any county-preservation
weight; the weight affects only the boundary placement at the county's
edges.
The default weight handles this correctly by treating county boundaries as
strong preferences rather than hard constraints.

\subsection{Component 4: The VRA Overlay}

For states where the Gingles three-prong test is met for a protected minority
group, the \dia{} mandates a separate VRA overlay run using VRASection \citep{b14paper}.

\paragraph{When the overlay applies.}
The overlay applies when:
\begin{enumerate}
  \item A minority group constitutes more than 50\% of the voting-age population
    in a geographic area that could form a compact congressional district
    (Gingles Prong~1: geographic compactness and size);
  \item The minority group votes cohesively (Gingles Prong~2: minority cohesion),
    as evidenced by precinct-level election results \citep{fekrazad2021}; and
  \item The majority votes sufficiently as a bloc to defeat the minority's
    preferred candidate in the absence of a majority-minority district
    (Gingles Prong~3: majority bloc voting).
\end{enumerate}

The first prong is evaluated algorithmically using the TIGER adjacency graph:
the Election Assistance Commission determines whether any set of adjacent
census tracts with combined minority VAP above 50\% can be assembled into
a geographically compact district.
The second and third prongs require precinct-level election data and are
evaluated by the Election Assistance Commission using ecological inference
\citep{imai2020} on the most recent federal election results.

\paragraph{How VRASection works.}
When the overlay applies, VRASection modifies the ratio-selection step of
the \hh{} bisection.
Normally, \ar{} selects the bisection ratio that minimizes the normalized
edge-cut $\mathrm{EC}_\text{norm}(i:j)$ at each level of the tree.
VRASection introduces a geographic alignment score:
\[
  A = \left|\text{MVAP}_\text{frac}(L) - \text{MVAP}_\text{frac}(R)\right|
\]
where $\text{MVAP}_\text{frac}(L)$ and $\text{MVAP}_\text{frac}(R)$ are
the minority VAP fractions in the left and right sub-regions.
The modified score function is:
\[
  \text{score}(\text{ratio}) = \mathrm{EC}_\text{norm} - w_\text{vra} \cdot A \cdot \max(\mathrm{EC}_\text{norm}, 1)
\]
where $w_\text{vra} = 0.40$ is calibrated below the \shaw{}--\miller{}
racial-predominance threshold: the alignment bonus reduces the compactness
score by at most 40\%, ensuring that compactness remains the dominant criterion.

\paragraph{Callais-mutex enforcement.}
The \callais{} preflight gate (Component 5, below) enforces that VRASection
and ProportionalSection---the partisan-signal mode---cannot be run in the
same redistricting process.
If a state triggers the VRA overlay, the partisan-proportionality mechanism
is disabled for that cycle.
This is not a policy choice; it is the architectural implementation of the
\callais{} disentanglement requirement.

\paragraph{Empirical validation.}
VRASection has been validated against Alabama (2020 census data).
The algorithm produces two majority-minority districts with Black VAP of
55.2\% and 52.8\%---satisfying the \allenmilligan{} requirement of two
Black-opportunity districts \citep{b14paper}.
The VRA weight $w_\text{vra} = 0.40$ is calibrated to produce this result
without exceeding the racial-predominance threshold.

\subsection{Component 5: The Callais Preflight Gate}

The \callais{} disentanglement requirement (at p.~36 of the majority opinion)
mandates that race-conscious and partisan signals cannot be combined in a
single redistricting run.
The \dia{} implements this as a hard computational gate---the
\texttt{callais\_preflight}---that enforces mutual exclusion at the CLI level.

If VRASection is activated (VRA overlay applies), ProportionalSection is
automatically disabled.
If ProportionalSection is activated (a state petitions for a proportionality
adjustment), VRASection is automatically disabled.
The two modes are mutually exclusive and the gate is enforced by the
\texttt{redist} binary at runtime: attempting to configure both will cause
the run to fail with a \texttt{[CONFIG]} error message referencing the
\callais{} disentanglement requirement.

This enforcement is documented in the plan manifest: the manifest records
which mode was used, and the absence of the other mode is cryptographically
verified by the manifest's hash chain.
A court reviewing a plan's compliance with \callais{} can verify compliance
by reading the manifest.

\subsection{Component 6: The GCD-Incompatibility Fallback}

Not all seat counts have factorizations that map cleanly onto a \hh{} bisection
tree.
\citet{b13paper} identifies that states where $\gcd(k, k-1) > 1$---where
the seat count and the seat count minus one share a common factor---present
factorization-spine difficulties.

More practically: states with large prime seat counts create bisection trees
with one very large leaf.
Texas ($k = 38 = 2 \times 19$) requires a first-level 2-way split followed
by a 19-way partition of each half.
The 19-way partition is a direct $k$-way METIS run, which is computationally
and geographically equivalent to the unweighted multi-way algorithm.
For Texas, the \hh{} tree is dominated by a single 19-way leaf, and the
first-level 2-way split is the only structural contribution of the \ar{} approach.

States with $k$ equal to a prime number are the limiting case.
If $k$ is prime, the only factorization is $1 \times k$: the entire state
is a single ``region'' that must be partitioned directly into $k$ districts
without any hierarchical bisection.
For such states, \ar{} reduces to a direct $k$-way METIS partition with
population balance tolerance $\texttt{ubvec}[0] = 1.001$.

\citet{b13paper} identifies the following states where the current seat count
is prime or where the largest prime factor is unusually large:
Wyoming ($k=1$, trivial), Alaska ($k=1$, trivial), Vermont ($k=1$, trivial),
Delaware ($k=1$, trivial), North Dakota ($k=1$, trivial), South Dakota ($k=1$,
trivial), Montana ($k=2$), Maine ($k=2$), New Hampshire ($k=2$),
Rhode Island ($k=2$), Hawaii ($k=2$), Idaho ($k=2$), West Virginia ($k=2$),
Nebraska ($k=3$), New Mexico ($k=3$), Kansas ($k=4 = 2 \times 2$, clean),
Nevada ($k=4$, clean), Utah ($k=4$, clean), Arkansas ($k=4$, clean).

For $k = 1$ states, redistricting is trivial: the entire state is one district.
For $k = 2$ states, the \hh{} tree requires a single 2-way split: the algorithm
bisects the state into two equal-population halves.
For $k = 3$ and larger primes, the fallback applies.

\paragraph{The prime-seat-count fallback.}
For states where $k$ is prime and $k \geq 3$, the \dia{} specifies the
following fallback:
\begin{enumerate}
  \item Run a direct $k$-way METIS partition with \texttt{ubvec}[0] = 1.001.
  \item Use edge weights from B.2 (geographic boundary lengths) and the
    county-preservation default ($\alpha_\text{county} = 2.0$).
  \item Use the content-derived seed.
\end{enumerate}

This fallback preserves all the properties of the \ar{} default
(population equality, compactness, county preservation, zero political inputs)
without the hierarchical bisection structure that the prime factorization
provides.
The loss is the multi-chamber compatibility property (the factorization spine):
a prime-seat-count state cannot nest its congressional districts cleanly
within a state-legislative factorization hierarchy.
This is a secondary concern; it does not affect the constitutional validity
of the resulting map.

The fallback applies to at most 6 states in the current apportionment
(those with $k$ prime and $k \geq 3$).
Nebraska ($k=3$) and New Mexico ($k=3$) are the most significant cases.
Congress could address the factorization incompatibility in these states
by providing that a state with a prime seat count may apportion by either
the fallback or a two-level split into $\lceil k/2 \rceil$ and $\lfloor k/2 \rfloor$
sub-regions, with the split selection made by the Election Assistance Commission
to maximize geographic compactness.

\subsection{The Full Statutory Text}

The complete proposed text of the \dia{} follows.
The operative provision is the one-sentence statutory command in subsection (c).
The remaining provisions are definitions and administrative details necessary
for implementation.

\begin{statute}
\noindent\textbf{Districting Integrity Act}\\
\textit{Proposed 2 U.S.C.\ \usca{}\ 2a(c)--(g) [new].}

\medskip
\noindent\textbf{(c) Intrastate Apportionment Method.}
Each State shall apportion its congressional districts by applying the
Huntington-Hill priority rule to its census tracts, using as input only
the prime factorization of the State's apportioned seat count and the
federal TIGER/Line adjacency graph for the census year in which the
apportionment was performed, with a seed determined by the formula
\texttt{SHA-256(census\_release\_id || `DIA\_SEED\_V1')}.

\medskip
\noindent\textbf{(d) County Preservation Default.}
In applying subsection (c), the State shall use a county boundary weight
of $\alpha_\text{county} = 2.0$, as specified in regulations promulgated by
the Election Assistance Commission, unless the Election Assistance Commission
has approved a different weight pursuant to subsection (f).

\medskip
\noindent\textbf{(e) VRA Overlay.}
In any State where the Election Assistance Commission determines that the
three-prong test of \textit{Thornburg v.\ Gingles}, 478 U.S. 30 (1986),
is met for a protected minority group, the State shall apply the
VRASection overlay to the allocation of bisection ratios, using a geographic
alignment weight of $w_\text{vra} = 0.40$, prior to running the algorithm
specified in subsection (c).
A State to which this subsection applies shall not simultaneously apply any
partisan-proportionality adjustment.

\medskip
\noindent\textbf{(f) Variance Petitions.}
A State may petition the Election Assistance Commission for a variance from
the county preservation weight specified in subsection (d) upon demonstrating
that the default weight is inconsistent with the State's constitution or
with subsection (e).
The Commission shall grant or deny the petition within 30 days.

\medskip
\noindent\textbf{(g) Prime Factorization Fallback.}
If the State's apportioned seat count $k$ is a prime number greater than 2,
the State shall apply a direct $k$-way METIS partition with population balance
tolerance \texttt{ubvec}[0] = 1.001, using the edge weights, county preservation
default, and seed formula specified in subsections (c), (d), and the
content-derived seed of subsection (c).

\medskip
\noindent\textbf{(h) Plan Verification.}
Each State shall submit its redistricting plan to the Election Assistance
Commission within 30 days of the release of the decennial census redistricting
data.
The plan submission shall include the plan manifest specified by the Election
Assistance Commission, which shall record all parameters, the seed, the TIGER
graph version, and the resulting district assignments, in a format that permits
independent verification of compliance with subsection (c).
\end{statute}

\subsection{The Decade-Stable Map Property}

\citet{b15paper} establishes that 67\% of states produce \ar{} maps that are
census-stable: running the algorithm on the 2000, 2010, and 2020 census data
produces the same district boundaries (to within a single-tract boundary
variation at the margins) in all three census years.
This property means that, in two-thirds of states, the map drawn after the
2030 census will closely resemble the map that would have been drawn after
the 2020 census, providing a decade of stable representation.

The 33\% of states that are not census-stable are those experiencing
significant population redistribution between census decades---primarily
Sun Belt states (Texas, Florida, Arizona, Nevada) with high population
growth and internal migration.
In these states, the \dia{} map will change when the new census data shifts
the factorization of the seat count (due to reapportionment) or shifts the
population distribution enough to change which bisection ratios are optimal.
This is not a bug.
It is the intended behavior: the map should update when the population changes
significantly, and the algorithm should produce the correct map for the
current census data without regard to what the previous map looked like.

The census-stability property is a feature of the algorithm, not a statutory
requirement.
The \dia{} does not mandate stability; it mandates the correct algorithm.
Stability is the empirical consequence of applying the correct algorithm
to a state with a stable population distribution.
